\section*{अध्याय \ref{ch:g4:magnets-and-compass} --- चुंबक और दिक्सूचक}

\begin{solution}{exo:g4:magnets-and-compass:1}
\omterm{def:g4:magnets-and-compass:poles}{ध्रुव} \omterm{def:g1:magnets:magnet}{चुंबक} के दोनों ताक़तवर सिरे हैं: \omterm{def:g4:magnets-and-compass:poles}{उत्तरी ध्रुव} (N, अकसर लाल)
और \omterm{def:g4:magnets-and-compass:poles}{दक्षिणी ध्रुव} (S)। क्लिप ध्रुवों पर --- यानी सिरों पर --- गुच्छा
बनाती हैं, बीच में नहीं।
\end{solution}

\begin{solution}{exo:g4:magnets-and-compass:2}
विपरीत \omterm{def:g4:magnets-and-compass:poles}{ध्रुव} \omterm{def:g1:magnets:magnet}{आकर्षित} करते हैं; एक जैसे \omterm{def:g4:magnets-and-compass:poles}{ध्रुव} दूर धकेलते हैं। चिपकने
का मतलब है कि विपरीत \omterm{def:g4:magnets-and-compass:poles}{ध्रुव} आमने-सामने थे: N के सामने S।
\end{solution}

\begin{solution}{exo:g4:magnets-and-compass:3}
N के सामने N: वे दूर धकेल देते हैं (एक जैसे \omterm{def:g4:magnets-and-compass:poles}{ध्रुव} दूर धकेलते हैं)।
N के सामने S: वे खिंचकर चिपक जाते हैं।
\end{solution}

\begin{solution}{exo:g4:magnets-and-compass:4}
क्योंकि एक विशाल \omterm{def:g1:magnets:magnet}{चुंबक} उसके ध्रुवों को क़तार में खींच रहा है: ख़ुद
पृथ्वी एक \omterm{def:g1:magnets:magnet}{चुंबक} है, जिसके चुंबकीय \omterm{def:g4:magnets-and-compass:poles}{ध्रुव} गोले के ऊपरी और निचले सिरों के
पास हैं।
\end{solution}

\begin{solution}{exo:g4:magnets-and-compass:5}
उत्तर, पूर्व, दक्षिण, पश्चिम। उत्तर की ओर मुँह हो तो दाईं ओर पूर्व
है।
\end{solution}

\begin{solution}{exo:g4:magnets-and-compass:6}
सुई ख़ुद एक नन्हा \omterm{def:g1:magnets:magnet}{चुंबक} है और वह \emph{सबसे पास} वाले तेज़ खिंचाव की
बात मानती है। पास रखा कोई \omterm{def:g1:magnets:magnet}{चुंबक} या लोहे का बड़ा \omterm{def:g3:measuring-mass:mass}{द्रव्यमान} दूर बैठे
पृथ्वी-चुंबक को मात दे देता है, और सुई उत्तर के बारे में झूठ बोलने
लगती है।
\end{solution}

\begin{solution}{exo:g4:magnets-and-compass:7}
सुई पर \omterm{def:g1:magnets:magnet}{चुंबक} का एक \omterm{def:g4:magnets-and-compass:poles}{ध्रुव} बीस बार, हमेशा एक ही दिशा में रगड़ो, और हर
रगड़ के बीच \omterm{def:g1:magnets:magnet}{चुंबक} उठा लो --- वह नन्हा \omterm{def:g1:magnets:magnet}{चुंबक} बन जाती है। फिर उसे पानी
के कटोरे में तैरते कॉर्क के टुकड़े पर रख दो: तैरता कॉर्क पानी से बनी
\omterm{def:g3:levers-and-scales:lever}{धुरी} है, और सुई घूमकर उत्तर--दक्षिण की ओर ठहर जाती है।
\end{solution}

\begin{solution}{exo:g4:magnets-and-compass:8}
वह उत्तर की ओर इशारा करती है। जिस \omterm{def:g4:magnets-and-compass:compass}{दिक्सूचक} की सुई दोपहर वाली \omterm{def:g1:light-and-shadows:shadow}{छाया} से
मेल खाती है, वह सच बोल रहा है --- और धूप वाले दिन अकेली \omterm{def:g1:light-and-shadows:shadow}{छाया} भी
मोटे तौर पर \omterm{def:g4:magnets-and-compass:compass}{दिक्सूचक} का काम कर सकती है।
\end{solution}

\begin{solution}{exo:g4:magnets-and-compass:9}
वह कोहरे में, बादलों के नीचे, चाँद-रहित रातों में और हर तट से दूर भी
काम करता है --- सुई को आकाश से कुछ नहीं चाहिए। नाविक आख़िरकार हर
मौसम में खुले महासागर के पार अपनी राह पकड़े रह सके।
\end{solution}

\begin{solution}{exo:g4:magnets-and-compass:10}
पास से साथी की बेल्ट का लोहे का बकसुआ पृथ्वी-चुंबक को मात दे देता है:
सुई उत्तर की जगह बकसुए की ओर इशारा करती है। उस पर भरोसा करके वह भटक
जाएगी --- झूठा ``उत्तर'' उसकी बाईं ओर पड़ता है, इसलिए वह अपनी राह
बाईं ओर, यानी पश्चिम की ओर मोड़ती चली जाएगी। हल: \omterm{def:g4:magnets-and-compass:compass}{दिक्सूचक} बकसुए से
दूर हटकर (कुछ क़दम अलग जाकर) पढ़ो, फिर चलो।
\end{solution}

\begin{solution}{exo:g4:magnets-and-compass:11}
रगड़ी हुई छोटी सुई के \emph{दो} \omterm{def:g4:magnets-and-compass:poles}{ध्रुव} होते हैं --- उसके सिरे N और S
की तरह बरताव करते हैं (ऐसी दो सुइयाँ तैराकर देखो, वे सिरे से सिरा
जोड़कर चिपक जाती हैं, विपरीत सिरे आपस में \omterm{def:g1:magnets:magnet}{आकर्षित} होते हुए)। यानी
\omterm{def:g1:magnets:magnet}{चुंबक} छोटा कर देने से कोई \omterm{def:g4:magnets-and-compass:poles}{ध्रुव} हटता नहीं, और उसे बीच से काटने पर
दो पूरे छोटे \omterm{def:g1:magnets:magnet}{चुंबक} बन जाते हैं: काटकर अकेला \omterm{def:g4:magnets-and-compass:poles}{ध्रुव} ढूँढ़ने की तलाश कभी
ख़त्म नहीं होती --- हर कटाई तुम्हें दोनों \omterm{def:g4:magnets-and-compass:poles}{ध्रुव} फिर से थमा देती है।
\end{solution}
